% Scalar representations and an exact hierarchy for the full free tree.
% The beta-Mellin identity is restored from the earlier archived derivation
% discrete_eit_transfer_target.tex, not asserted as a new discovery.
% The explicit finite-z tail below is supplied here for convenient use.
\section{A scalar beta series retaining the full free tree}
\label{hfk:section}

Start the original free friend tree from an edge, with no inert stub.
Let $C_n$ be its number of nonleaves at population $n$, and put
$a_n=\mathbb E C_n$. Thus $a_2=0$, $a_3=1$, and $a_4=4/3$.
Every $a_n$ is rational. More explicitly, if $\nu_n(T)$ is the law on
labelled trees grown from the initial edge, then
\[
 \nu_{n+1}(T+w)\mathrel{+}=\nu_n(T)\frac1n
       \sum_{u\sim_Tw}\frac1{d_T(u)},\qquad
 a_n=\sum_T\nu_n(T)C(T).
\]
Here $T+w$ denotes attaching the new labelled vertex $n+1$ to the
existing vertex $w$.
One may instead merge isomorphic trees, retaining all transition
multiplicities. This is a finite exact recurrence at every population;
it is not a closed recurrence of fixed order for $a_n$.

Define the ordinary scalar series
\[
 G(z)=\sum_{n\ge2}(n-1)a_nz^{n-2},\qquad |z|<1.
\]
This free-tree generating function is distinct from the uniform
newborn envelope used to define $H(T)$ in the computation theorem.
The Yule population process is independent of the embedded tree chain,
because its holding times at population $n$ are independent exponentials
of rate $n$. With $z=1-e^{-t}$ this gives
\begin{equation}\label{hfk:yule}
 F_2(t):=\mathbb E_{\rm edge}C(t)
   =(1-z)^2G(z).
\end{equation}
Indeed $\Pr(N(t)=n)=(n-1)e^{-2t}(1-e^{-t})^{n-2}$.
Consequently, by Tonelli, for real $s>0$ including divergent values,
\begin{equation}\label{hfk:mellin}
 \mathcal M(s):=\int_0^\infty e^{-st}F_2(t)\,dt
  =\sum_{n\ge2}(n-1)a_n\,\mathrm B(s+2,n-1)
  =\Gamma(s+2)\sum_{n\ge2}
       a_n\frac{\Gamma(n)}{\Gamma(n+s+1)}.
\end{equation}
The established mean exponents\cite{Long,Uzay},
$F_2(t)=\exp(\mu t+o(t))$ and $a_n=n^{\mu+o(1)}$, therefore give
\begin{equation}\label{hfk:abscissa}
 \mu=\inf\{s>0:\mathcal M(s)<\infty\}.
\end{equation}
The series converges absolutely and defines a holomorphic function on
$\Re s>\mu$, and diverges for real $0<s<\mu$. Section~\ref{cr:section}
also proves divergence at $s=\mu$. Equivalently,
\[
 \mu=\lim_{z\uparrow1}
       \frac{\log G(z)}{-\log(1-z)}-2.
\]

The exact scalar organization \eqref{hfk:mellin} was already derived in
the project's earlier exploratory work. It is recorded here to connect
that representation to the current pressure theorem. It is a full-model
series with universal beta factors, but its coefficients retain all tree
dependence. No hypergeometric term ratio, holonomic recurrence, or finite
special-function equation for these coefficients has been proved.
In particular \eqref{hfk:abscissa} is an abscissa characterization, not an
equation obtained by setting $\mathcal M(\mu)$ to a finite constant.

\begin{proposition}[Explicit geometric tail at a fixed interior argument]
\label{hfk:fixed-z}
Let $0<z<1$, let $N\ge2$, and put $M=N-1$. If
\[
 G_N(z)=\sum_{n=2}^{N}(n-1)a_nz^{n-2},
\]
then
\begin{align*}
 0\le G(z)-G_N(z)
 \le z^M\left[
  \frac{M(M+1)}{1-z}
  +\frac{(2M+1)z}{(1-z)^2}
  +\frac{z(1+z)}{(1-z)^3}
 \right].
\end{align*}
Multiplying the same interval by $(1-z)^2$ encloses
$F_2(-\log(1-z))$. At rational $z$, the partial sum and both error
endpoints are rational.
\end{proposition}
\begin{proof}
Every finite tree has at least two leaves, so $0\le a_n\le n-2$.
Writing $m=n-2$, the omitted sum is bounded by
$\sum_{m\ge M}m(m+1)z^m$. Substitute $m=M+j$ and use
$\sum z^j=(1-z)^{-1}$,
$\sum jz^j=z(1-z)^{-2}$, and
$\sum j^2z^j=z(1+z)(1-z)^{-3}$.
\end{proof}

For fixed $z<1$ the tail is $O_z(N^2z^N)$, a geometric convergence
bound. Its deterioration as $z\uparrow1$ is essential: this estimate
does not provide a convergence rate for the limiting logarithmic
singularity that defines $\mu$. Similarly, the beta-series summand is
$n^{\mu-s-1+o(1)}$, so its unaccelerated convergence can be arbitrarily
slow as $s\downarrow\mu$.

\subsection{An exact difference and differential hierarchy}

The scalar coefficients admit a useful hierarchy beyond the full-state
probability recurrence. For a finite free tree $T$, put
\[
 h_T(w)=\sum_{u\sim_Tw}\frac1{d_T(u)},\qquad
 (\mathcal L f)(T)=\sum_{w\in T}h_T(w)[f(T+w)-f(T)].
\]
Here $\sum_wh_T(w)=|T|$, and $C(T)$ counts vertices of degree at
least two. Let $f_r=\mathcal L^r C$ and
$u_{r,n}=\mathbb E f_r(T_n)$, so $u_{0,n}=a_n$.

\begin{proposition}[Full-model recurrence hierarchy]
\label{hfk:hierarchy}
For every $r\ge0$ and $n\ge2$,
\begin{equation}\label{hfk:moment-recurrence}
 u_{r,n+1}=u_{r,n}+\frac1n u_{r+1,n}.
\end{equation}
All its coefficients and initial values $u_{r,2}$ are rational and
are computable from finitely many tree states. If
\[
 G_r(z)=\sum_{n\ge2}(n-1)u_{r,n}z^{n-2},
\]
then these series converge for $|z|<1$ and satisfy
\begin{equation}\label{hfk:differential-hierarchy}
 (1-z)G_r'(z)-2G_r(z)=G_{r+1}(z),\qquad r\ge0.
\end{equation}
In particular $G_0=G$ is the full-tree series above. There is also
the exact finite operator-polynomial formula
\begin{equation}\label{hfk:operator-factorial}
 a_n=\frac{[(\mathcal L+2)_{n-2}C](T_2)}{(n-1)!},
 \qquad
 G(z)=\sum_{m\ge0}\frac{[(\mathcal L+2)_mC](T_2)}{m!}z^m.
\end{equation}
The rising factorial in this formula is a finite polynomial in the
full-tree generator, not a scalar hypergeometric parameter.
\end{proposition}
\begin{proof}
Conditional on $T_n=T$, the next attachment has probability
$h_T(w)/n$ at $w$. Thus for every observable $f$,
\[
 \mathbb E[f(T_{n+1})\mid T_n]
 =f(T_n)+\frac1n(\mathcal Lf)(T_n).
\]
Apply this to $f_r$ and average. Each $f_r$ at a fixed finite tree
uses at most $r$ further births, hence only finitely many rational
transitions. Induction also gives the coarse bound
$|f_r(T)|\le2^r(|T|)_{r+1}$. Therefore $G_r$ and its derivative
converge absolutely inside the unit disc. Comparing coefficients
using~\eqref{hfk:moment-recurrence} proves
\eqref{hfk:differential-hierarchy}.

On states of population $k$, the operator $I+\mathcal L/k$ is the
one-step transition operator. Composing these operators for
$k=2,\ldots,n-1$ and evaluating at the edge gives $a_n$.
They are polynomials in the same generator and commute algebraically;
their product is $(\mathcal L+2)_{n-2}/(n-1)!$. Every finite
composition is well defined on the finitely many states it visits.
This proves~\eqref{hfk:operator-factorial} without assuming a
bounded-operator functional calculus or an entire semigroup expansion.
\end{proof}

The first two levels make the environmental dependence explicit.
Write $\ell_v$ for the number of leaf neighbours of a nonleaf vertex
$v$, and set $\Lambda=\mathcal LC$. At the initial edge,
$\Lambda(T_2)=2$. At every tree with at least three vertices,
\[
 \Lambda(T)=\sum_{v\text{ nonleaf}}\frac{\ell_v}{d_v},\qquad
 K(T)=\sum_{v\text{ nonleaf}}\frac{\ell_v}{d_v^2},\qquad
 J(T)=\sum_{v\text{ nonleaf}}
          \frac{h_T(v)(d_v-\ell_v)}{d_v(d_v+1)}.
\]
An exact calculation gives
\begin{equation}\label{hfk:second-level}
 (\mathcal L\Lambda)(T)=\tfrac12\Lambda(T)-K(T)+J(T).
\end{equation}
Indeed, attachment to a nonleaf $v$ changes its contribution to
$\Lambda$ by $(d_v-\ell_v)/[d_v(d_v+1)]$ and occurs at rate
$h_T(v)$. Promotion of a leaf whose neighbour is $v$ changes
$\Lambda$ by $1/2-1/d_v$ and has aggregate rate $\ell_v/d_v$.
Adding these changes proves~\eqref{hfk:second-level}.

Consequently, if $b_n=\mathbb E\Lambda(T_n)$, then
\[
 n(a_{n+1}-a_n)=b_n\quad(n\ge2),\qquad
 n(b_{n+1}-b_n)=\tfrac12b_n-\mathbb EK(T_n)+\mathbb EJ(T_n)
 \quad(n\ge3).
\]
The edge requires the separate value
$(\mathcal L\Lambda)(T_2)=-2$; the displayed nonleaf formula is
not applied to it. These recurrences give
\[
 a_2=0,\quad a_3=1,\quad a_4=\frac43,\quad
 a_5=\frac{19}{12},\quad a_6=\frac{647}{360}.
\]

There is no closure in the core count alone. For example, take two
adjacent nonleaf vertices with four attached leaves. Allocations
$(3,1)$ and $(2,2)$ both give six-vertex trees with $C=2$, but their
$\Lambda$ values are $5/4$ and $4/3$, respectively. Their next-step
core-increase probabilities are therefore $5/24$ and $2/9$.
This rules out a statewise transition law determined by $C$ alone.
It does not rule out a different finite recurrence for edge-start
expectations. The hierarchy above is exact; no finite scalar closure
has been established.
Section~\ref{fgo:section} strengthens this obstruction by proving that
no finite invariant linear observable space containing $C$ exists,
while preserving the distinction from a special edge-start scalar relation.

Finally, for $n\ge3$ put $\beta_n=b_n/a_n$. Since
$0\le\Lambda(T)\le C(T)$ in these states, $0\le\beta_n\le1$.
The first recurrence and $a_3=1$ give
\begin{equation}\label{hfk:logarithmic-average}
 a_n=\prod_{k=3}^{n-1}\left(1+\frac{\beta_k}{k}\right),
 \qquad
 \mu=\lim_{n\to\infty}\frac1{\log n}
                  \sum_{k=3}^{n-1}\frac{\beta_k}{k}.
\end{equation}
For the limit, use the established $a_n=n^{\mu+o(1)}$ and
$0\le x-\log(1+x)\le x^2/2$ for $x\ge0$; the accumulated
logarithmic error is bounded by $\tfrac12\sum_{k\ge3}k^{-2}$.
No pointwise limit $\beta_n\to\mu$, or rate for this logarithmic
average, is asserted.
